\begin{corollary}[Concentration around the convex fitted vector]
Under the assumptions of (PR2), let $b_*\in\mathbb R^n$ be the
training vector of a minimizer of the full empirical path objective.
This vector is unique, and
\[
 \mathbb E_\Pi\frac1n\|A_\theta c-b_*\|_2^2
 \le C\lambda\nu\log\frac e{\lambda\nu}. \tag{PR4}
\]
The same bound holds for the training vector of the ordinary
posterior mean.
\end{corollary}
\begin{proof}
Let $p(b)$ be the minimum signed-unit-ridge mass representing the
training vector $b$, allowing an arbitrary finite number of atoms.
Its domain is a linear subspace of $\mathbb R^n$ and $p$ is a norm
on that subspace. The objective
$\|b-Y\|_2^2/n+2\lambda p(b)$ is strictly convex there, so its
minimizing training vector is unique. The norm subgradient at its
minimum and the exact quadratic expansion give, for every feasible $b$,
\[
 \frac1n\|b-Y\|_2^2+2\lambda p(b)-J_*
 \ge\frac1n\|b-b_*\|_2^2.
\]
Every amplitude implementation has $P(c)\ge p(A_\theta c)$.
Apply (PR2) and then Jensen's inequality for the posterior mean.
\end{proof}
